\documentclass[../main]{subfiles}
\begin{document}

\chapter{Exposición Integrador Máquinas Térmicas}
\img{images/16/caldera.jpeg}{Calderas}


    \begin{table}[h!]
        \centering
        \begin{tabular}{|p{5cm} | p{10cm} |} \hline
            Tema & Subtema \\ \hline \hline 
             Seguridad & Accesorios de seguridad en calderas.\\ \hline 
             Calderas de Uso industrial & Estructura  \\ \hline
             Condensadores y Quemadores & combustión, chimenea y caja de humos, intercambiadores de calor\\ \hline
             Máquinas Frigoríficas& Ciclo inverso de Carnot\\ \hline
             Centrales térmicas y equipamiento auxiliar para calderas& precalentadores, desgasificadores y bombas.\\ \hline 
        \end{tabular}
        \caption{Temas de Calderas}
        
    \end{table}




\section{Contenido de la presentación.}
\begin{enumerate}
    \item \emph{Principio de funcionamiento:} Explicar cada una de las partes que componen a la Máquina Térmica, tanto la parte de la alimentación que entrega la energía, con especificaciones técnicas y un ejemplo concreto. 
    \item Ciclo de funcionamiento: Explicar como funciona como circula el calor y en que se convierte, o el princípio físico que caracteriza su funcionamiento.
    \item \emph{Fotos}: Buscar al menos de 10 fotos de este tipo de máquinas.
    \item \emph{Historia:} Historia de su fabricación, primeras aplicaciones en la industria, inventor creador, y primeros uso de este tipo de Máquina Térmica.
    \item \emph{Características Técnicas}: consumo de energía potencia, energía combustible, combustión, presión, caudal de funcionamiento, presiones,etc.
    \item \emph{Mantenimiento}: Explicar como se realiza el mantenimiento de esta Máquina Térmica.
    \item \emph{Ventajas y desventajas:} Analizar las ventajas y desventajas del uso de este tipo de equipo. 
    \item Mercado actual:
    \begin{itemize}
        \item ¿Dónde se compra?
        \item ¿Cuál es el precio?
        \item ¿Que características tiene cada diferentes  modelo comercial?
    \end{itemize}
    \item \emph{Conclusiones:} Resumen con conclusiones generales después de la presentación de la clase.
    \item \emph{Preguntas:} Se responderán las preguntas que los alumnos tengan con respecto a la presentación.
    \item Subir archivo pdf con el informe y powerpoint (PPT), a classroom con el nombre de los integrantes.
\end{enumerate}


\end{document}




